\clearpage
\phantomsection
\chapter*{题目索引}
\addcontentsline{toc}{chapter}{题目索引}
\par\medskip\noindent\textbf{第1章}\par
\BookIndexEntry{co-ex1-1}{例 1-1　硬件直接执行哪一种程序 【2015】}
\BookIndexEntry{co-ex1-2}{例 1-2　主频相同，怎样由指令组成求时间 【2022】}
\BookIndexEntry{co-ex1-3}{例 1-3　CPU 加快后，I/O 时间怎样处理 【2012】}
\BookIndexEntry{co-self1}{自测}
\par\medskip\noindent\textbf{第2章}\par
\BookIndexEntry{co-ex2-1}{例 2-1　整数与小数分别换基数 〖2.1-001〗}
\BookIndexEntry{co-ex2-2}{例 2-2　无符号短整数扩展 【2012】}
\BookIndexEntry{co-ex2-3}{例 2-3　负的 short 转成 unsigned int 【2025】}
\BookIndexEntry{co-ex2-4}{例 2-4　减法器的两路输入不直接送负数符号 〖2.2-018〗}
\BookIndexEntry{co-ex2-5}{例 2-5　十六进制补码相减 【2018】}
\BookIndexEntry{co-ex2-6}{例 2-6　有符号减法为何仍会产生借位 【2023】}
\BookIndexEntry{co-ex2-7}{例 2-7　负小数怎样组成单精度机器数 【2022】}
\BookIndexEntry{co-ex2-8}{例 2-8　同为负浮点数时怎样比较 【2014】}
\BookIndexEntry{co-ex2-9}{例 2-9　哪些步骤可能造成指数越界 【2015】}
\BookIndexEntry{co-self2}{自测}
\par\medskip\noindent\textbf{第3章}\par
\BookIndexEntry{co-ex3-1}{例 3-1　DRAM 的复用地址怎样计数 【2014】}
\BookIndexEntry{co-ex3-2}{例 3-2　已安装容量不是全部可编址空间 【2011】}
\BookIndexEntry{co-ex3-3}{例 3-3　同一模块不等于立即发生冲突 【2015】}
\BookIndexEntry{co-ex3-4}{例 3-4　ROM 与 RAM 分区要分别计算 【2009】}
\BookIndexEntry{co-ex3-5}{例 3-5　7200转每分不是7200转每秒 【2015】}
\BookIndexEntry{co-ex3-6}{例 3-6　读后写的数组访问怎样计缺失率 【2016】}
\BookIndexEntry{co-ex3-7}{例 3-7　8路组相联需要几个比较器 【2022】}
\BookIndexEntry{co-ex3-8}{例 3-8　Cache 总容量还包括哪些位 【2015】}
\BookIndexEntry{co-ex3-9}{例 3-9　页号的十六进制不能直接当表中十进制 【2022】}
\BookIndexEntry{co-ex3-10}{例 3-10　TLB 标记不包含页内偏移 【2024】}
\BookIndexEntry{co-ex3-11}{例 3-11　交换内外层循环，命中率一定降低吗 【2023】改编}
\BookIndexEntry{co-self3}{自测}
\par\medskip\noindent\textbf{第4章}\par
\BookIndexEntry{co-ex4-1}{例 4-1　扩展操作码应逐层扣除哪些前缀 【2022】}
\BookIndexEntry{co-ex4-2}{例 4-2　相对转移用哪个 PC 【2009】}
\BookIndexEntry{co-ex4-3}{例 4-3　先确定结构布局，再看低字节 【2012】}
\BookIndexEntry{co-ex4-4}{例 4-4　RISC 设计能否支持流水线 【2025】}
\BookIndexEntry{co-ex4-5}{例 4-5　工具链顺序 【2022】}
\BookIndexEntry{co-self4}{自测}
\par\medskip\noindent\textbf{第5章}\par
\BookIndexEntry{co-ex5-1}{例 5-1　哪些元件负责运算，哪些元件保存状态 【2023】}
\BookIndexEntry{co-ex5-2}{例 5-2　从单总线图确定控制信号与取指时序 【2022】}
\BookIndexEntry{co-ex5-3}{例 5-3　间接操作数经过哪些部件 【2019】}
\BookIndexEntry{co-ex5-4}{例 5-4　互斥类为什么还要一个空操作编码 【2012】}
\BookIndexEntry{co-ex5-5}{例 5-5　控制存储器地址字段怎样估位宽 【2014】}
\BookIndexEntry{co-ex5-6}{例 5-6　缺页后为什么重试当前指令 【2021】}
\BookIndexEntry{co-self5}{自测}
\par\medskip\noindent\textbf{第6章}\par
\BookIndexEntry{co-ex6-1}{例 6-1　最慢的一段决定时钟下限 【2009】}
\BookIndexEntry{co-ex6-2}{例 6-2　装入后立即使用需要一个气泡 〖5.6-013〗}
\BookIndexEntry{co-ex6-3}{例 6-3　同一段序列里同时有数据与控制等待 【2023】}
\BookIndexEntry{co-ex6-4}{例 6-4　硬件多线程不等于多核 【2022】}
\BookIndexEntry{co-self6}{自测}
\par\medskip\noindent\textbf{第7章}\par
\BookIndexEntry{co-ex7-1}{例 7-1　一个总线周期占两个时钟周期 【2009】}
\BookIndexEntry{co-ex7-2}{例 7-2　已经给出 MT/s 就不再乘4 【2025】}
\BookIndexEntry{co-ex7-3}{例 7-3　不支持突发，就要反复准备每个数据 【2023】}
\BookIndexEntry{co-ex7-4}{例 7-4　峰值带宽与整次事务的有效速率 【2024】}
\BookIndexEntry{co-self7}{自测}
\par\medskip\noindent\textbf{第8章}\par
\BookIndexEntry{co-ex8-1}{例 8-1　命令字走哪组线 【2012】}
\BookIndexEntry{co-ex8-2}{例 8-2　定时查询占多少 CPU 时间 【2011】}
\BookIndexEntry{co-ex8-3}{例 8-3　响应顺序与完成顺序为什么可以不同}
\BookIndexEntry{co-ex8-4}{例 8-4　允许延迟不是额外的服务开销 【2014】}
\BookIndexEntry{co-ex8-5}{例 8-5　周期挪用不等于整块期间封住主存 【2020】}
\BookIndexEntry{co-ex8-6}{例 8-6　DMA 的准备与完成仍需要软件 【2019】}
\BookIndexEntry{co-self8}{自测}
\clearpage
\phantomsection
\chapter*{方法索引}
\addcontentsline{toc}{chapter}{方法索引}
\par\medskip\noindent\textbf{第1章}\par
\BookIndexEntry{co-01-01}{从程序看计算机的功能部件}
\BookIndexEntry{co-01-02}{存储程序怎样变成一次次指令执行}
\BookIndexEntry{co-01-03}{性能计算：先确定谁的时间在改变}
\par\medskip\noindent\textbf{第2章}\par
\BookIndexEntry{co-02-01}{位串与数值之间还差一个解释规则}
\BookIndexEntry{co-02-02}{定点编码：范围先于运算规则}
\BookIndexEntry{co-02-03}{截断、扩展与 C 类型转换}
\BookIndexEntry{co-02-04}{加法器怎样同时支持加与减}
\BookIndexEntry{co-02-05}{CF 与 OF 分别判断什么}
\BookIndexEntry{co-02-06}{乘除运算必须保留的中间量}
\BookIndexEntry{co-02-07}{IEEE 754：字段既规定精度，也规定特殊情况}
\BookIndexEntry{co-02-08}{浮点加减：对阶丢位以后还要判断舍入}
\par\medskip\noindent\textbf{第3章}\par
\BookIndexEntry{co-03-01}{为什么一个程序同时使用几层存储器}
\BookIndexEntry{co-03-02}{SRAM、DRAM 与 Flash 保存信息的方式}
\BookIndexEntry{co-03-03}{地址引脚不总等于地址位数}
\BookIndexEntry{co-c03-h35}{MAR 容量由地址空间决定}
\BookIndexEntry{co-03-04}{单体多字和多体交叉怎样提高带宽}
\BookIndexEntry{co-03-05}{主存芯片怎样连到 CPU}
\BookIndexEntry{co-03-06}{磁盘与 SSD 的一次访问花在哪些时间上}
\BookIndexEntry{co-03-07}{Cache 命中率怎样变成时间}
\BookIndexEntry{co-03-08}{地址拆分决定块能放在哪里}
\BookIndexEntry{co-03-09}{LRU 更新的是访问次序，不只是装入次序}
\BookIndexEntry{co-03-10}{写命中与写缺失是两个独立问题}
\BookIndexEntry{co-03-11}{虚拟存储器怎样让程序看到地址空间}
\BookIndexEntry{co-03-12}{页表与 TLB：地址偏移为何原样保留}
\BookIndexEntry{co-03-13}{段式与段页式地址为什么需要界限检查}
\BookIndexEntry{co-memory-synthesis}{跨页数组访问：页、Cache块与遍历次序分别计算}
\par\medskip\noindent\textbf{第4章}\par
\BookIndexEntry{co-04-01}{指令字段怎样决定可表示的操作和地址}
\BookIndexEntry{co-04-02}{从形式地址到有效地址}
\BookIndexEntry{co-04-03}{大端、小端和对齐解决不同问题}
\BookIndexEntry{co-04-04}{RISC 与 CISC：设计倾向不代替具体实现}
\BookIndexEntry{co-04-05}{编译、汇编、链接怎样确定程序中的地址}
\BookIndexEntry{co-04-06}{分支和循环怎样改变 PC}
\BookIndexEntry{co-04-07}{过程调用为什么要保存返回位置}
\par\medskip\noindent\textbf{第5章}\par
\BookIndexEntry{co-05-01}{寄存器的名字说明它暂存什么}
\BookIndexEntry{co-05-02}{指令周期、机器周期与时钟周期}
\BookIndexEntry{co-05-03}{一条指令怎样变为微操作}
\BookIndexEntry{co-05-04}{控制信号由组合逻辑还是微程序产生}
\BookIndexEntry{co-05-05}{异常、中断与返回位置}
\par\medskip\noindent\textbf{第6章}\par
\BookIndexEntry{co-06-01}{延迟和完成间隔是两个量}
\BookIndexEntry{co-06-02}{阶段寄存器保存的不只是一份数据}
\BookIndexEntry{co-06-03}{冒险先按冲突对象分类}
\BookIndexEntry{co-06-04}{超标量与动态流水线允许哪些重叠}
\BookIndexEntry{co-06-05}{一条指令处理多个数据，还是多个指令流}
\BookIndexEntry{co-06-06}{硬件线程、多核与共享内存}
\par\medskip\noindent\textbf{第7章}\par
\BookIndexEntry{co-07-01}{位宽、频率、传送率分别给出什么}
\BookIndexEntry{co-07-02}{一次总线事务从地址到应答}
\BookIndexEntry{co-c07-h75}{如何读总线时序图}
\par\medskip\noindent\textbf{第8章}\par
\BookIndexEntry{co-08-01}{端口是接口中可访问的位置}
\BookIndexEntry{co-08-02}{查询方式：等待和搬运都由 CPU 安排}
\BookIndexEntry{co-08-03}{中断方式：先请求，再响应，再运行服务程序}
\BookIndexEntry{co-08-04}{DMA：数据不经过 CPU 寄存器逐次搬运}
